% Compile with: latexmk -lualatex -pdf praat-pipeline.tex
\documentclass[11pt,a4paper]{article}

% ---------- Fonts & Languages ----------
\usepackage{fontspec}
% \setmainfont{Charis SIL}

\usepackage{polyglossia}
\setmainlanguage{english}

% ---------- Linguistics ----------
\usepackage{tipauni}

% ---------- Graphics & drawing ----------
\usepackage{graphicx}
\usepackage{tikz}
\usetikzlibrary{calc}
\usepackage{xcolor}

% ---------- Math ----------
\usepackage{amsmath}

% ---------- Layout ----------
\usepackage{geometry}
\geometry{margin=2.5cm}
\usepackage{hyperref}
\hypersetup{colorlinks=true, linkcolor=blue, urlcolor=blue}

% ---------- Custom colors for the tiers ----------
\definecolor{phoncolor}{RGB}{230, 240, 255}
\definecolor{wordcolor}{RGB}{240, 255, 230}

% ---------- Reusable macro for a phoneme-tier box ----------
% Usage: \phonbox{startX}{endX}{label}
% Centres the label robustly instead of relying on `node[midway]`
% on a rectangle, which behaves differently across TikZ versions.
\newcommand{\phonbox}[3]{%
    \draw[fill=phoncolor, draw=black] (#1,-0.5) rectangle (#2,0);
    \node at ({(#1+#2)/2},-0.25) {#3};
}

\title{Praat $\to$ LuaLaTeX Pipeline \\
       \large Spectrograms with TikZ-drawn alignment tiers}
\author{Your Name}
\date{\today}

\begin{document}
\maketitle

This worked example shows the pipeline used by phoneticians and computational
linguists to publish aligned acoustic data: export a spectrogram from Praat as
a PDF, embed it with \texttt{graphicx}, and draw the phoneme / word / gloss
tiers over it with TikZ.

\section{Step 1 --- Export from Praat}

In Praat, select the Sound and the TextGrid together, open \textbf{View \&
Edit}, show the spectrogram and the tiers you want, then \textbf{File $\to$
Save as PDF\ldots}  Call the file \texttt{spectrogram.pdf}.

For automated workflows, export the TextGrid as a text or CSV file and read it
from Python with \texttt{parselmouth} or \texttt{textgridtools}. That gives you
the exact start/end time of every segment in seconds, which you can feed
directly into the TikZ coordinates below.

\section{Step 2 --- Embed and overlay in LaTeX}

The example below uses a gray placeholder box in place of a real spectrogram
so this file compiles out of the box. Replace it with
\verb|\includegraphics| once you have a real \texttt{spectrogram.pdf}.

\begin{center}
\begin{tikzpicture}[x=1cm, y=1cm]

    % 1. THE SPECTROGRAM
    % Real version:
    % \node[anchor=south west, inner sep=0] (image) at (0,0)
    %   {\includegraphics[width=12cm]{spectrogram.pdf}};
    %
    % Placeholder (delete when using a real image):
    \draw[fill=gray!10, draw=gray] (0,0) rectangle (12,4);
    \node at (6,2) {\textbf{[Insert Praat Spectrogram PDF Here]}};

    % 2. THE ALIGNMENT TIERS (drawn below the spectrogram)

    % --- PHONEME TIER ---
    \phonbox{0}{1.5}{k}
    \phonbox{1.5}{3}{\textipa{\ae}}
    \phonbox{3}{4.5}{t}
    \node[anchor=east] at (-0.2,-0.25) {\textbf{Phonemes:}};

    % --- WORD TIER ---
    \draw[fill=wordcolor, draw=black] (0,-1.2) rectangle (4.5,-0.7);
    \node at (2.25,-0.95) {cat};
    \node[anchor=east] at (-0.2,-0.95) {\textbf{Words:}};

    % --- GLOSS TIER ---
    \draw[draw=black] (0,-1.9) rectangle (4.5,-1.4);
    \node at (2.25,-1.65) {\textit{`cat'}};
    \node[anchor=east] at (-0.2,-1.65) {\textbf{Gloss:}};

    % 3. TIME STAMPS
    \draw[thick] (0,0) -- (0,-2.1);
    \draw[thick] (4.5,0) -- (4.5,-2.1);
    \node[below] at (0,-2.1) {0.00\,s};
    \node[below] at (4.5,-2.1) {0.45\,s};

\end{tikzpicture}
\end{center}

\section{Step 3 --- Scaling the coordinates}

TikZ draws in abstract units, but the spectrogram sits in seconds. If your
audio file is $T$ seconds long and you embedded the spectrogram at width $W$
cm, then an event starting at time $t$ seconds has $x$-coordinate
\[
x \;=\; \frac{t}{T} \cdot W \quad \text{(cm)}.
\]
Example: a vowel starting at $0.3$~s in a $1.5$~s file, with a $12$~cm wide
spectrogram, starts at $x = (0.3/1.5) \times 12 = 2.4$~cm.

For anything more than a handful of segments, let Python do the arithmetic.
A short script can read the TextGrid and spit out
\verb|\phonbox{...}{...}{...}| lines directly, which you then
\verb|\input{}| into the TikZ picture. LaTeX does the drawing; Python does
the math.

\section{Related packages}

\begin{itemize}
    \item \texttt{tikz-dependency} --- for dependency parsing: draws labeled
          arcs over text (useful for spaCy or UDPipe output).
    \item \texttt{expex}, \texttt{linguex} --- alternatives to \texttt{gb4e}
          for numbered glosses; some departments prefer them.
    \item \texttt{forest} --- phrase-structure trees (already in the main kit).
\end{itemize}

\end{document}
